% !TEX program = lualatex
% Stat 150 — chapter-by-chapter notes scaffold
% ============================================================
%  Stat 150 — Chapter-by-Chapter Reading Notes
%  Durrett, \emph{Essentials of Stochastic Processes} (3rd ed.)
%  Keshav Ramamurthy — Fall 2026
% ------------------------------------------------------------
%  HOW TO USE (one growing document, one banner per chapter)
%
%    Find the chapter banner and the section heading below it,
%    then type your notes directly underneath:
%
%        Big idea: every Cauchy sequence converges iff ...
%
%    Structure the important ideas with the note environments
%    (numbered within the current chapter):
%
%        \begin{definition} ... \end{definition}
%        \begin{theorem} ... \end{theorem}
%        \begin{example}  ... \end{example}
%        \begin{remark}   ... \end{remark}
%        \begin{idea}     ... \end{idea}
%
%    Add a missing section:  \notesec{§1.4}{Title}
%    Add a chapter:          \chapterbanner{1}{Chapter 1}{Title}{description}
% ============================================================
\documentclass[11pt]{article}

\usepackage{fontspec}
\usepackage{microtype}
\usepackage{mathtools,amssymb,amsfonts,amsthm,bm}
\usepackage{enumitem}
\usepackage{booktabs,array,xcolor}
\usepackage{geometry,fancyhdr,setspace}
\usepackage{etoolbox}
\usepackage[most]{tcolorbox}
\usepackage[hidelinks]{hyperref}
\usepackage[nameinlink,noabbrev]{cleveref}

% ---------- Page and paragraph rhythm (matches lecture/solutions templates) ----------
\geometry{margin=1in,headheight=15pt}
\setstretch{1.12}
\setlength{\parindent}{0pt}
\setlength{\parskip}{0.55em plus 0.12em minus 0.08em}
\setlength{\abovedisplayskip}{0.8em plus 0.2em minus 0.1em}
\setlength{\belowdisplayskip}{0.8em plus 0.2em minus 0.1em}
\allowdisplaybreaks[3]
\setlist{leftmargin=*,topsep=0.35em,itemsep=0.15em,parsep=0pt}

% ---------- Palette (matches reference cheatsheet) ----------
\definecolor{ink}{HTML}{1E293B}
\definecolor{navy}{HTML}{183B56}
\definecolor{paperblue}{HTML}{E8F1F8}
\definecolor{gold}{HTML}{FFF3CD}
\definecolor{mint}{HTML}{E8F5E9}
\definecolor{muted}{HTML}{52606D}
\definecolor{rule}{HTML}{C9D5DF}
\color{ink}

% ---------- Metadata ----------
\newcommand{\studentname}{Keshav Ramamurthy}
\newcommand{\coursename}{Stat 150}
\newcommand{\courselongname}{Stochastic Processes}
\newcommand{\textbookname}{Durrett, \emph{Essentials of Stochastic Processes} (3rd ed.)}
\newcommand{\term}{Fall 2026}

% ---------- Core notation (same as ~/.config/nvim/textemplate.tex) ----------
\newcommand{\N}{\mathbb{N}}
\newcommand{\Z}{\mathbb{Z}}
\newcommand{\Q}{\mathbb{Q}}
\newcommand{\R}{\mathbb{R}}
\newcommand{\C}{\mathbb{C}}
\newcommand{\F}{\mathbb{F}}
\newcommand{\K}{\mathbb{K}}
\newcommand{\E}{\mathbb{E}}
\newcommand{\Prob}{\mathbb{P}}
\newcommand{\1}{\mathbf{1}}
\newcommand{\eps}{\varepsilon}
\newcommand{\dd}{\,\mathrm{d}}
\newcommand{\abs}[1]{\left\lvert#1\right\rvert}
\newcommand{\norm}[1]{\left\lVert#1\right\rVert}
\newcommand{\ip}[2]{\left\langle#1,#2\right\rangle}
\newcommand{\set}[1]{\left\{#1\right\}}
\newcommand{\ceil}[1]{\left\lceil#1\right\rceil}
\newcommand{\floor}[1]{\left\lfloor#1\right\rfloor}
\newcommand{\given}{\,\middle\vert\,}
\newcommand{\indep}{\perp\!\!\!\perp}
\newcommand{\wh}[1]{\widehat{#1}}
\newcommand{\deriv}[2]{\frac{\mathrm{d}#1}{\mathrm{d}#2}}
\newcommand{\pderiv}[2]{\frac{\partial#1}{\partial#2}}
\newcommand{\lto}{\longrightarrow}
\newcommand{\st}{\text{ such that }}
\DeclareMathOperator{\Aut}{Aut}
\DeclareMathOperator{\End}{End}
\DeclareMathOperator{\Gal}{Gal}
\DeclareMathOperator{\Hom}{Hom}
\DeclareMathOperator{\Ker}{ker}
\DeclareMathOperator{\im}{im}
\DeclareMathOperator{\rank}{rank}
\DeclareMathOperator{\nullity}{nullity}
\DeclareMathOperator{\Span}{span}
\DeclareMathOperator{\tr}{tr}
\DeclareMathOperator{\diag}{diag}
\DeclareMathOperator{\spec}{spec}
\DeclareMathOperator{\ord}{ord}
\DeclareMathOperator{\supp}{supp}
\DeclareMathOperator{\diam}{diam}
\DeclareMathOperator{\Var}{Var}
\DeclareMathOperator{\Cov}{Cov}
\DeclareMathOperator*{\argmin}{arg\,min}
\DeclareMathOperator*{\argmax}{arg\,max}

% ---------- Note environments (numbered by chapter) ----------
\newcounter{chap}
\newtheoremstyle{notestyle}{0.7em}{0.7em}{\itshape}{}{}{\bfseries}{.5em}{}
\theoremstyle{notestyle}
\newtheorem{theorem}{Theorem}[chap]
\newtheorem{lemma}[theorem]{Lemma}
\newtheorem{proposition}[theorem]{Proposition}
\newtheorem{corollary}[theorem]{Corollary}
\theoremstyle{definition}
\newtheorem{definition}[theorem]{Definition}
\newtheorem{example}[theorem]{Example}
\newtheorem{remark}[theorem]{Remark}

% ---------- Note machinery ----------
% Chapter banner: \chapterbanner{1}{Chapter 1}{Groups and Subgroups}{Fraleigh: ...}
\newtcolorbox{chapterbox}[2]{%
  enhanced,
  colback=paperblue, colframe=navy, boxrule=0.7pt, arc=2.5pt,
  left=10pt, right=10pt, top=12pt, bottom=8pt,
  boxed title style={colback=navy, colframe=navy, arc=2pt, boxrule=0pt,
    left=8pt, right=8pt, top=3pt, bottom=3pt},
  attach boxed title to top left={yshift=-2.4mm, xshift=5mm},
  title={\sffamily\bfseries\color{white}\large #1\quad #2}}
\newcommand{\chapterbanner}[4]{%
  \clearpage
  \setcounter{chap}{#1}%
  \setcounter{theorem}{0}%
  \phantomsection
  \addcontentsline{toc}{section}{#2\quad #3}%
  \begin{chapterbox}{#2}{#3}
    {\small\sffamily\color{muted}#4\par}
  \end{chapterbox}
  \medskip\nopagebreak}

% Section heading: \notesec{§1.4}{Title}  →  navy heading + thin rule
\newcommand{\notesec}[2]{%
  \par\goodbreak\medskip\nopagebreak
  {\sffamily\bfseries\color{navy}#1\quad #2}%
  \par\nobreak\vspace{1.5pt}
  {\color{rule}\hrule height 0.6pt}\par\nobreak\smallskip\nopagebreak
  \phantomsection
  \addcontentsline{toc}{subsection}{#1\quad #2}}

% Quick structure inside notes
\newenvironment{claim}{\par\smallskip\noindent\textit{Claim.}\ }{\par\smallskip}
\newenvironment{idea}{\par\smallskip\noindent\textit{Idea.}\ }{\par\smallskip}
\newenvironment{proofsketch}{\par\smallskip\noindent\textit{Proof sketch.}\ }{\par\smallskip}

% ---------- Header ----------
\pagestyle{fancy}
\fancyhf{}
\fancyhead[L]{\coursename}
\fancyhead[C]{\courselongname\ — chapter notes}
\fancyhead[R]{\studentname}
\fancyfoot[C]{\thepage}

\begin{document}

% ---------- Title ----------
\begin{center}
  {\Large\sffamily\bfseries\color{navy}\coursename\ -- \courselongname}\\[3pt]
  {\sffamily\color{muted} Chapter-by-chapter notes on the important ideas}\\[3pt]
  {\small\sffamily\color{muted}\textbookname\quad\textbullet\quad \term\quad\textbullet\quad \studentname}
\end{center}
\medskip
{\color{rule}\hrule height 0.8pt}
\bigskip

{\small
\begin{tcolorbox}[colback=paperblue, colframe=rule, boxrule=0.4pt, arc=2pt,
  left=7pt, right=7pt, top=5pt, bottom=5pt,
  title={\sffamily\bfseries\color{navy} Taking notes},
  colbacktitle=paperblue]
\sffamily Find the chapter banner, then type under the section heading as you read:
\begin{center}
\texttt{\textbackslash notesec\{1A\}\{Span and Linear Independence\}}
\end{center}
Record big ideas with \texttt{idea}, definitions with \texttt{definition}, theorems with
\texttt{theorem}, \texttt{lemma}, \texttt{proposition}, examples with \texttt{example}, and
warnings with \texttt{remark} --- all numbered within the chapter.
Add a missing section by copying a \texttt{\textbackslash notesec} slot;
add a chapter with \texttt{\textbackslash chapterbanner}.
\end{tcolorbox}
}

\medskip
\tableofcontents

\chapterbanner{1}{Chapter 1}{Markov Chains}{Durrett:\ transition probabilities, classification of states, stationary distributions, limit behavior, exit distributions and times}

\notesec{§1.1}{Definitions and Examples}
% Big idea:
% Key definitions \& theorems:
% Examples / tricks:
% Questions:

\notesec{§1.2}{Multistep Transition Probabilities}
% Big idea:
% Key definitions \& theorems:
% Examples / tricks:
% Questions:

\notesec{§1.3}{Classification of States}
% Big idea:
% Key definitions \& theorems:
% Examples / tricks:
% Questions:

\notesec{§1.4}{Stationary Distributions}
% Big idea:
% Key definitions \& theorems:
% Examples / tricks:
% Questions:

\notesec{§1.5}{Detailed Balance Condition}
% Big idea:
% Key definitions \& theorems:
% Examples / tricks:
% Questions:

\notesec{§1.6}{Limit Behavior}
% Big idea:
% Key definitions \& theorems:
% Examples / tricks:
% Questions:

\notesec{§1.7}{Returns to a Fixed State}
% Big idea:
% Key definitions \& theorems:
% Examples / tricks:
% Questions:

\notesec{§1.8}{Proof of the Convergence Theorem*}
% Big idea:
% Key definitions \& theorems:
% Examples / tricks:
% Questions:

\notesec{§1.9}{Exit Distributions}
% Big idea:
% Key definitions \& theorems:
% Examples / tricks:
% Questions:

\notesec{§1.10}{Exit Times}
% Big idea:
% Key definitions \& theorems:
% Examples / tricks:
% Questions:

\notesec{§1.11}{Infinite State Spaces*}
% Big idea:
% Key definitions \& theorems:
% Examples / tricks:
% Questions:

\notesec{§1.12}{Chapter Summary}
% Big idea:
% Key definitions \& theorems:
% Examples / tricks:
% Questions:

\chapterbanner{2}{Chapter 2}{Poisson Processes}{Durrett:\ the exponential distribution, defining and constructing the Poisson process, compound Poisson, thinning and superposition}

\notesec{§2.1}{Exponential Distribution}
% Big idea:
% Key definitions \& theorems:
% Examples / tricks:
% Questions:

\notesec{§2.2}{Defining the Poisson Process}
% Big idea:
% Key definitions \& theorems:
% Examples / tricks:
% Questions:

\notesec{§2.3}{Compound Poisson Processes}
% Big idea:
% Key definitions \& theorems:
% Examples / tricks:
% Questions:

\notesec{§2.4}{Transformations}
% Big idea:
% Key definitions \& theorems:
% Examples / tricks:
% Questions:

\notesec{§2.5}{Chapter Summary}
% Big idea:
% Key definitions \& theorems:
% Examples / tricks:
% Questions:

\chapterbanner{3}{Chapter 3}{Renewal Processes}{Durrett:\ laws of large numbers, queueing theory applications, age and residual life}

\notesec{§3.1}{Laws of Large Numbers}
% Big idea:
% Key definitions \& theorems:
% Examples / tricks:
% Questions:

\notesec{§3.2}{Applications to Queueing Theory}
% Big idea:
% Key definitions \& theorems:
% Examples / tricks:
% Questions:

\notesec{§3.3}{Age and Residual Life*}
% Big idea:
% Key definitions \& theorems:
% Examples / tricks:
% Questions:

\notesec{§3.4}{Chapter Summary}
% Big idea:
% Key definitions \& theorems:
% Examples / tricks:
% Questions:

\chapterbanner{4}{Chapter 4}{Continuous Time Markov Chains}{Durrett:\ transition probabilities, limiting behavior, exit distributions and times, Markovian queues, networks}

\notesec{§4.1}{Definitions and Examples}
% Big idea:
% Key definitions \& theorems:
% Examples / tricks:
% Questions:

\notesec{§4.2}{Computing the Transition Probability}
% Big idea:
% Key definitions \& theorems:
% Examples / tricks:
% Questions:

\notesec{§4.3}{Limiting Behavior}
% Big idea:
% Key definitions \& theorems:
% Examples / tricks:
% Questions:

\notesec{§4.4}{Exit Distributions and Exit Times}
% Big idea:
% Key definitions \& theorems:
% Examples / tricks:
% Questions:

\notesec{§4.5}{Markovian Queues}
% Big idea:
% Key definitions \& theorems:
% Examples / tricks:
% Questions:

\notesec{§4.6}{Queueing Networks*}
% Big idea:
% Key definitions \& theorems:
% Examples / tricks:
% Questions:

\notesec{§4.7}{Chapter Summary}
% Big idea:
% Key definitions \& theorems:
% Examples / tricks:
% Questions:

\chapterbanner{5}{Chapter 5}{Martingales}{Durrett:\ conditional expectation, examples, gambling strategies and stopping times, applications}

\notesec{§5.1}{Conditional Expectation}
% Big idea:
% Key definitions \& theorems:
% Examples / tricks:
% Questions:

\notesec{§5.2}{Examples}
% Big idea:
% Key definitions \& theorems:
% Examples / tricks:
% Questions:

\notesec{§5.3}{Gambling Strategies, Stopping Times}
% Big idea:
% Key definitions \& theorems:
% Examples / tricks:
% Questions:

\notesec{§5.4}{Applications}
% Big idea:
% Key definitions \& theorems:
% Examples / tricks:
% Questions:

\chapterbanner{6}{Chapter 6}{Mathematical Finance}{Durrett:\ the binomial model, American options, the Black--Scholes formula, calls and puts}

\notesec{§6.1}{Two Simple Examples}
% Big idea:
% Key definitions \& theorems:
% Examples / tricks:
% Questions:

\notesec{§6.2}{Binomial Model}
% Big idea:
% Key definitions \& theorems:
% Examples / tricks:
% Questions:

\notesec{§6.3}{Concrete Examples}
% Big idea:
% Key definitions \& theorems:
% Examples / tricks:
% Questions:

\notesec{§6.4}{American Options}
% Big idea:
% Key definitions \& theorems:
% Examples / tricks:
% Questions:

\notesec{§6.5}{Black--Scholes Formula}
% Big idea:
% Key definitions \& theorems:
% Examples / tricks:
% Questions:

\notesec{§6.6}{Calls and Puts}
% Big idea:
% Key definitions \& theorems:
% Examples / tricks:
% Questions:

\chapterbanner{7}{Appendix A}{Review of Probability}{Durrett:\ probability spaces, independence, random variables and distributions, expectation, limit theorems}

\notesec{§A.1}{Probabilities, Independence}
% Big idea:
% Key definitions \& theorems:
% Examples / tricks:
% Questions:

\notesec{§A.2}{Random Variables, Distributions}
% Big idea:
% Key definitions \& theorems:
% Examples / tricks:
% Questions:

\notesec{§A.3}{Expected Value, Moments}
% Big idea:
% Key definitions \& theorems:
% Examples / tricks:
% Questions:

\notesec{§A.4}{Integration to the Limit}
% Big idea:
% Key definitions \& theorems:
% Examples / tricks:
% Questions:

\end{document}
